\chapter{相关理论与关键技术}

为了使后续方法设计建立在清晰、可追溯的技术基础之上，本章从工程图纸检索任务所涉及的几个核心方向出发，对本文所依赖的理论与关键技术进行梳理。具体包括 CLIP 视觉语义表示、参数高效微调、向量检索与重排、局部注意机制、OCR 字段级多模态融合以及 Grad-CAM 可解释性分析。通过对这些技术的归纳，可以明确本文方法设计的理论来源、工程可行性与改进空间。

\section{工程图纸检索任务特点}

工程图纸检索与自然图像检索存在显著差异。自然图像通常以纹理、颜色和语义对象为主要判别线索，而工程图纸更强调轮廓、投影视图、尺寸标注、标题栏字段和图面布局等结构性信息。因此，在工程图纸场景中，单纯依赖自然图像任务中常见的全局语义表征，往往难以充分区分外形相近但结构关系不同的图纸样本。

此外，工程图纸还具有三个典型特点。第一，图纸的类别边界往往由局部几何关系、剖视方式和投影组合决定，判别线索更细粒度。第二，标题栏中包含图号、零件名称、材料和比例等强语义信息，这些文本字段对排序具有补充价值。第三，图纸在来源转换过程中容易出现黑底、边框、截屏痕迹和模板差异，这些噪声会干扰视觉编码与文本提取。正是由于上述特点，工程图纸检索更适合采用“视觉召回 + 结构增强 + 文本辅助”的综合策略，而不是简单套用通用图像检索流程。

\section{CLIP 模型原理}

CLIP 通过大规模图文对对比学习，将图像与文本映射到统一嵌入空间中~\citep{radford2021clip}。对于输入图像 $I$ 和文本 $T$，图像编码器与文本编码器分别输出特征向量 $\bm{v}$ 和 $\bm{t}$，再通过余弦相似度进行匹配：
\begin{equation}
  s(I,T)=\frac{\bm{v}^{\top}\bm{t}}{\|\bm{v}\|_2\|\bm{t}\|_2}.
\end{equation}
该机制使 CLIP 具备较强的跨任务泛化能力，因此可以作为工程图纸检索的视觉主干。但对于结构主导型工程图纸，仅依赖全局语义表示仍难以充分刻画局部线稿结构特征。

\section{参数高效微调与 Adapter}

为了在保留大模型通用知识的同时适应特定领域，参数高效微调技术应运而生。CLIP-Adapter 通过在视觉或文本分支后附加轻量 MLP 实现快速领域适配~\citep{gao2024clipadapter}；CoOp 则通过可学习提示优化文本侧语义表达~\citep{zhou2022coop}。这类方法降低了全量微调的成本，也为垂直工业场景提供了更稳定的迁移策略。

在本文系统的前期实现中，曾参考“视觉主干 + 轻量适配”的思路，但在后续检索优化过程中，研究重点逐步转向结构建模与多模态重排。其主要原因在于，相较于重新训练更复杂的模型，在既有检索框架中引入结构信息与字段级语义信息，更有利于获得可解释、可验证且可复现的性能改进。

\section{工程图纸与 CAD 检索相关研究}

围绕工程图纸、草图和 CAD 数据的检索研究，近年来逐渐形成了两条较有代表性的思路。第一类工作关注跨模态或跨视图匹配，例如从二维正投影视图检索三维 CAD 模型、从草图检索图像或模型等。这类研究表明，预训练视觉语义模型与图结构描述在工程数据中具有良好的迁移潜力~\citep{sain2023clipforallthings,mahajan2025orthocad,heidari2025geometriccad}。第二类工作则更加关注内容检索流程本身，包括图像表示、相似度建模、候选重排与索引组织等~\citep{zhang2023cbirsurvey,liu2024matchingreview,han2024crossmodalreview}。

不过，已有研究大多聚焦公开草图数据集、自然图像检索任务或 CAD 模型对齐任务，而对工业企业内部常见的工程图纸数据关注较少。实际工程图纸往往具有模板差异大、转换噪声多、标题栏布局复杂、结构相似样本密集等特点，这使得现有方法在直接迁移时面临较强的领域鸿沟。本文工作正是在这样的背景下展开：一方面利用预训练模型微调后的视觉表征作为统一起点，另一方面结合结构重排序与字段级 OCR 辅助，面向真实工程图纸场景构建可运行的检索系统。

\section{向量检索与重排技术}

向量检索将输入图纸编码为高维嵌入向量，再通过近似最近邻搜索完成候选召回。分层小世界图是目前常用的近似最近邻索引结构之一，兼顾检索精度与效率~\citep{malkov2018hnsw}。向量数据库则在索引之上进一步提供元数据过滤、持久化和服务化能力~\citep{pan2024vectordbsurvey,qdrant2024site}。

然而，初始召回主要依赖全局特征，因此往往需要在候选阶段引入结构重排、多模态排序修正等机制，以提高前排结果质量。对于工程图纸任务，这一步尤其重要，因为许多零件图在全局轮廓上高度接近，仅靠单一视觉向量容易出现前排混淆。

在工程实现层面，向量检索系统不仅要追求召回质量，还要兼顾构建效率、更新成本和在线响应时间。面向大规模向量库的 GPU 加速索引和相似度计算研究表明，合理的近似搜索结构与批量编码方式能够显著降低系统延迟~\citep{johnson2019faiss}。因此，本文在算法设计之外，也将 GPU 化建库、离线缓存和候选重排机制纳入系统实现考虑范围。

\section{局部注意机制}

注意力机制最初用于序列建模任务，随后在视觉 Transformer 中被广泛用于图像 patch token 之间的关系建模~\citep{vaswani2017attention}。视觉 Transformer 将图像划分为 patch token 后进行全局建模，使模型能够在较大范围内捕获局部区域之间的依赖关系。本文引入的掩码引导局部注意机制并非重新设计复杂的注意力网络，而是利用图纸预处理阶段生成的掩码，将结构保留区域映射到 patch 空间，在特征聚合时对主体结构区域赋予更高权重。与直接增加深层注意力模块相比，该方案具有更低的额外计算开销，也更便于在现有工程系统中部署。

\section{OCR 与标题栏字段级多模态融合}

工程图纸标题栏通常包含图号、零件名、材料、比例和数量等关键信息。这些字段往往比普通说明文字更具类别区分性。传统 OCR 加分方案通常直接比较整段文本相似度，容易受到模板说明、技术要求等噪声影响。本文将 OCR 从“通用文本加分”升级为“标题栏字段级辅助”，使多模态信号更加聚焦于强语义字段，而非被大段说明文本干扰。

\section{Grad-CAM 可解释性分析}

Grad-CAM 通过对目标分数进行梯度回传，生成神经网络在当前任务中的关注热力图~\citep{selvaraju2017gradcam}。本文以查询图与候选图的嵌入相似度作为目标分数，计算 patch 级 Grad-CAM 热图，用于分析模型在相似检索中的决策依据。该工具一方面可以为论文中的注意力机制提供直观证据，另一方面也能辅助误检案例分析与后续模型诊断。

\section{本章评述}

综合已有研究可以看出，预训练视觉模型为工程图纸检索提供了良好的基础表征能力，参数高效微调降低了领域适配成本，向量检索与候选重排为大规模图库提供了可行的工程路径，OCR 和可解释性分析则为结果增强与误检诊断提供了补充工具。但同时也应看到，现有研究中仍存在三个不足：其一，对“显式清洗”与“无清洗主体增强”两种策略的系统对比不充分；其二，结构信息往往被简单视为附加特征，而缺乏与视觉召回紧密结合的精排机制；其三，多模态文本信号在工业图纸中的使用常常停留在通用文本比对层面，缺少针对标题栏字段的针对性设计。

\section{本章小结}

本章围绕本文研究所需的关键理论与技术基础进行了系统梳理。总体而言，CLIP 为工程图纸检索提供了统一的视觉表征起点，参数高效微调为领域适配提供了低成本路径，向量检索与候选重排构成了系统的基本工程框架，局部注意机制与结构增强为解决图纸主体表达不足问题提供了方法依据，OCR 字段级融合为引入强语义信息提供了补充路径，Grad-CAM 则为后续结果解释与误检分析提供了分析工具。上述内容共同构成了本文方法设计与系统实现的理论基础，下一章将在此基础上给出系统总体设计。
